\documentclass[11pt]{article}

\usepackage[margin=1in]{geometry}
\usepackage{amsmath,amssymb,bm,mathtools}
\usepackage{booktabs}
\usepackage{enumitem}
\usepackage{xcolor}
\usepackage{hyperref}

\title{ScaleBridge Phase E.0 --- E0-3A\\
Generic RC-Network Mathematical Contract\\
\large Draft v0.1 for Scientific Lock}
\author{ScaleBridge / PhD Code Framework}
\date{2026-08-23}

\newcommand{\R}{\mathbb{R}}
\newcommand{\Vc}{\mathcal{V}_{C}}
\newcommand{\Vb}{\mathcal{V}_{B}}
\newcommand{\Er}{\mathcal{E}_{R}}
\newcommand{\Qset}{\mathcal{Q}}
\newcommand{\Zset}{\mathcal{Z}}
\newcommand{\Sc}{\mathcal{S}_{c}}
\newcommand{\Sr}{\mathcal{S}_{r}}
\newcommand{\diag}{\operatorname{diag}}

\begin{document}
\maketitle

\section*{Status and authority}

This document is the proposed generic mathematical contract for ScaleBridge
Phase E.0 E0-3.  It is intentionally method-neutral.

The controlled PINODE/EPSR RC mathematics is a scientific reference for:
\begin{itemize}
    \item convective/radiative heat semantics,
    \item observed versus latent thermal states,
    \item reciprocal resistance coupling,
    \item positive $R/C$ parameterization,
    \item zone and building total-energy balances,
    \item the separation of IND/DEP1/DEP2 information architectures.
\end{itemize}

The controlled paper's RestaurantFastFood/Buffalo/Dining/Kitchen dimensions,
fixed forcing vectors, fixed two-zone coupling, fixed $\Delta t=300$ s, and
paper-specific DEP2 allocation are \emph{not} part of this generic contract.

The authoritative data boundary is the Phase-D final product interpreted
through the E0-2 Phase-E canonical contract.  Phase B and Phase C remain the
upstream providers of physical measurements, aggregation semantics, and
heat-input models; E0-3 does not duplicate those upstream models.

\section*{1. Orthogonal axes: RC topology versus information architecture}

A central ScaleBridge rule is that the physical RC topology and the
Phase-D information architecture are orthogonal.

An RC topology specifies:
\begin{itemize}
    \item thermal state/capacitance nodes,
    \item prescribed boundary-temperature nodes,
    \item reciprocal resistance edges,
    \item heat-injection ports and routing,
    \item observed-state selection.
\end{itemize}

The information architecture specifies which measured/estimated signals are
available:
\begin{itemize}
    \item All-to-One,
    \item IND,
    \item DEP1,
    \item DEP2.
\end{itemize}

Therefore a ``2R2C DEP1'' model is not a separate hand-written equation family.
It is:
\[
\boxed{
\text{generic RC graph}
+
\text{2R2C graph template}
+
\text{DEP1 signal map}.
}
\]

Likewise, 1R1C, 2R2C, 3R2C, and 4R3C are named graph instantiations of the
same generic RC-network contract.

\section*{2. Generic thermal graph}

Define the thermal graph
\[
\boxed{
\mathcal{G}_{RC}
=
(\Vc,\Vb,\Er,\Zset,\zeta,\Qset,\Gamma,H).
}
\]

Here:
\begin{align*}
\Vc &= \text{set of capacitive/state nodes},\\
\Vb &= \text{set of prescribed boundary-temperature nodes},\\
\Er &= \text{set of reciprocal resistance edges},\\
\Zset &= \text{set of modeled zones},\\
\zeta(v) &\in \Zset \text{ is the zone membership of state node }v,\\
\Qset &= \text{set of named thermal-power ports},\\
\Gamma &= \text{heat-routing operator},\\
H &= \text{observation/measurement operator}.
\end{align*}

Each capacitive node $v\in\Vc$ has temperature $T_v(t)$ and capacitance
\[
\boxed{C_v>0.}
\]

Each resistance edge is undirected:
\[
e=\{p,q\}\in\Er,
\]
with
\[
\boxed{R_e>0,\qquad G_e=\frac{1}{R_e}>0.}
\]

A resistance edge may connect:
\begin{enumerate}[label=(\roman*)]
    \item two state nodes,
    \item a state node and a boundary node.
\end{enumerate}

Boundary nodes have prescribed temperatures and are not integrated as states.

\section*{3. Reciprocal conductive heat flow}

For an edge $e=\{p,q\}$, define the heat flow into $p$ from $q$ as
\[
\boxed{
\Phi_{q\rightarrow p}
=
\frac{T_q-T_p}{R_e}.
}
\]

The reciprocal heat flow is
\[
\Phi_{p\rightarrow q}
=
-\Phi_{q\rightarrow p}.
\]

Hence every internal resistance edge is exactly energy-conservative:
\[
\boxed{
\Phi_{q\rightarrow p}
+
\Phi_{p\rightarrow q}
=
0.
}
\]

For a cross-zone edge, one and only one resistance parameter belongs to the
undirected edge:
\[
\boxed{R_{ij}=R_{ji}.}
\]

The physical implementation must never create two independently trainable
parameters for the two directions of the same thermal connection.

\section*{4. Generic nodewise RC equation}

For every capacitive node $v\in\Vc$,
\[
\boxed{
C_v\dot T_v
=
\sum_{e=\{v,j\}\in\Er}
\frac{T_j-T_v}{R_e}
+
\sum_{s\in\Qset}
\Gamma_{vs}q_s.
}
\tag{E0.3-1}
\]

The first term contains conductive/resistive exchange.
The second term contains externally supplied thermal-power signals.

The heat-routing coefficient $\Gamma_{vs}$ is dimensionless.
It may be fixed or parameterized, but its scientific meaning must be explicit.

\section*{5. Matrix form}

Let
\[
x=
\begin{bmatrix}
T_{v_1}&\cdots&T_{v_n}
\end{bmatrix}^{\!\top},
\qquad n=|\Vc|,
\]
and
\[
C=\diag(C_{v_1},\ldots,C_{v_n}).
\]

Let $T_B$ collect prescribed boundary temperatures and let $q$ collect the
named thermal-power ports.

Construct the weighted graph Laplacian and partition it with respect to state
and boundary nodes:
\[
L=
\begin{bmatrix}
L_{CC}&L_{CB}\\
L_{BC}&L_{BB}
\end{bmatrix}.
\]

Then the complete continuous-time dynamics are
\[
\boxed{
C\dot x
=
-L_{CC}x
-L_{CB}T_B
+\Gamma q.
}
\tag{E0.3-2}
\]

Equivalently,
\[
\boxed{
\dot x
=
C^{-1}
\left(
-L_{CC}x
-L_{CB}T_B
+\Gamma q
\right).
}
\tag{E0.3-3}
\]

No numerical discretization is part of E0-3.  Euler, Heun/RK2, RK4, the
locked sixth-order method, and physical/normalized time handling belong to
E0-5.

\section*{6. Units and sign convention}

The locked units are:
\begin{align*}
T &:\ ^\circ\mathrm{C}\ \text{or K on a consistent temperature basis},\\
R &:\ \mathrm{K/W},\\
C &:\ \mathrm{J/K},\\
q &:\ \mathrm{W},\\
\dot T &:\ \mathrm{K/s}.
\end{align*}

Temperature differences in conductive terms are numerically identical in
degrees Celsius and kelvin.

A positive heat port injects thermal energy into its routed node.
For HVAC:
\[
\boxed{
Q_{AC}>0 \text{ means heating},
\qquad
Q_{AC}<0 \text{ means cooling}.
}
\]

$P_{HVAC}$ is electrical power and is never a term in
(E0.3-1)--(E0.3-3).

\section*{7. Upstream signal semantics}

Phase D, interpreted through the E0-2 contract, is the authoritative training
data source.

For zone $i$, the primary grouped Phase-D thermal signals are:
\begin{align*}
Q_{AC,i} &:\ \text{signed HVAC thermal delivery},\\
Q_{ZIC,i} &:\ \text{grouped convective internal gain},\\
Q_{ZIR,i} &:\ \text{grouped radiative internal gain},\\
Q_{Sol1,i} &:\ \text{convective solar contribution},\\
Q_{Sol2,i} &:\ \text{radiative solar contribution}.
\end{align*}

When those exact bindings exist, define
\[
\boxed{
Q_{c,i}
=
Q_{AC,i}+Q_{ZIC,i}+Q_{Sol1,i},
}
\tag{E0.3-4}
\]
and
\[
\boxed{
Q_{r,i}
=
Q_{ZIR,i}+Q_{Sol2,i}.
}
\tag{E0.3-5}
\]

If Phase D exposes component-level heat signals instead of grouped signals,
the heat-set membership in the Phase-D manifest determines the corresponding
sum.

Crucially, the generic equations sum only over signals that are present under
the authoritative E0-2 binding/availability contract.  A missing applicable
signal is not silently imputed as zero.

A structurally inapplicable/structurally-zero channel may contribute no port
only when that status is explicit in the upstream contract.

\section*{8. Heat routing is a separate physical contract}

The RC graph and the heat-routing operator $\Gamma$ are separate.

For a named source $s$, define its destination set
\[
\mathcal{D}(s)\subseteq\Vc.
\]

For a conservative thermal source,
\[
\boxed{
\Gamma_{vs}\ge 0,
\qquad
\sum_{v\in\mathcal{D}(s)}\Gamma_{vs}=1.
}
\tag{E0.3-6}
\]

This guarantees that the source is redistributed but not created or destroyed.

For learnable conservative routing, one valid parameterization is
\[
\boxed{
\Gamma_{vs}
=
\frac{\exp(\rho_{vs})}
{\sum_{u\in\mathcal{D}(s)}\exp(\rho_{us})}.
}
\tag{E0.3-7}
\]

Routing parameters are physical/model parameters, not generic Optuna
hyperparameters.

\subsection*{8.1 Convective/radiative default semantics}

The primary ScaleBridge convention is:
\begin{itemize}
    \item convective heat enters the zone-air node;
    \item radiative heat is routed according to the selected topology-specific
          routing contract.
\end{itemize}

For a one-capacitance reduced model, the conservative default is that all
radiative heat enters the single state.

An optional reduced-order effective radiative participation factor
\[
0\le \eta_r^{(1C)}\le1
\]
may be enabled explicitly, giving
\[
Q_{\mathrm{eff}}^{(1C)}
=
Q_c+\eta_r^{(1C)}Q_r.
\]
This option is an effective reduced-order approximation.  The main
energy-conservative 1C choice is
\[
\boxed{\eta_r^{(1C)}=1.}
\]

For a two-state air/mass routing profile,
\[
\boxed{
\Gamma_{a,r}=1-\eta_r,
\qquad
\Gamma_{m,r}=\eta_r,
\qquad
0\le\eta_r\le1.
}
\tag{E0.3-8}
\]

\section*{9. Observation and latent-state contract}

Each modeled zone $i$ has a distinguished zone-air state node $a(i)$.

Phase D supplies the observed zone-air temperature:
\[
\boxed{
y_i=T_{a(i)}.
}
\]

Collecting all modeled zones:
\[
\boxed{y=Hx.}
\tag{E0.3-9}
\]

Every additional capacitive node not selected by $H$ is a physical latent
thermal state.

Thus:
\begin{itemize}
    \item 1C: the full state may be observed;
    \item $n$C with $n>1$: non-air capacitance states are latent unless an
          upstream measurement contract explicitly says otherwise.
\end{itemize}

Phase D does not fabricate latent-state targets.
Initialization/estimation of latent states belongs to E0-4 and later
method-specific mathematical contracts.

\section*{10. 1C versus nC organization}

For each modeled zone $i$, let
\[
\Vc_i=\{v\in\Vc:\zeta(v)=i\}.
\]

Define
\[
n_{C,i}=|\Vc_i|.
\]

The organization is:
\begin{align*}
\text{1C} &: n_{C,i}=1,\\
\text{nC} &: n_{C,i}\ge 2.
\end{align*}

This distinction affects state observability and heat routing, but not the
generic node equation (E0.3-1).

\section*{11. Information architecture: All-to-One}

All-to-One is an upstream aggregation/information architecture, not a new RC
physics equation.

There is one modeled aggregate zone $A$:
\[
\Zset=\{A\}.
\]

The aggregate Phase-D product supplies:
\[
T_A,\quad T_o,\quad
\{q_{A,s}\}_{s\in\mathcal{S}_A}.
\]

Any RC graph template may be instantiated on this aggregate zone.

\section*{12. Information architecture: IND}

In IND, each Phase-D product models one zone independently.

For zone $i$:
\[
\Zset=\{i\},
\]
and there are no modeled inter-zone RC edges.

The forcing map is a named map, not a hard-coded tensor:
\[
\boxed{
v_i^{IND}(t)
=
\left\{
T_B(t),
q_{i,s}(t):
s\in\mathcal{S}_i
\right\}.
}
\tag{E0.3-10}
\]

Different zones may have different sets $\mathcal{S}_i$.

\section*{13. Information architecture: DEP1}

DEP1 models multiple zones jointly.

Let
\[
\Zset=\{1,\ldots,N_z\}.
\]

Every modeled zone contributes its observed state and all local Phase-D
controls/disturbances that are actually available.

The joint named forcing map is
\[
\boxed{
v^{DEP1}(t)
=
\left\{
T_B(t),
q_{i,s}(t):
i\in\Zset,\ s\in\mathcal{S}_i
\right\}.
}
\tag{E0.3-11}
\]

Physical inter-zone coupling is supplied by the RC graph, not by Phase D.

For an inter-zone edge connecting nodes $v\in i$ and $w\in j$:
\[
\boxed{
\Phi_{j\rightarrow i}
=
\frac{T_w-T_v}{R_{ij}},
\qquad
\Phi_{i\rightarrow j}
=
-\Phi_{j\rightarrow i}.
}
\tag{E0.3-12}
\]

There is one shared positive resistance for that undirected edge.

\section*{14. Information architecture: DEP2}

DEP2 uses the same physical RC graph class as DEP1.

What changes is the disturbance information map:
\begin{itemize}
    \item modeled-zone states remain zone-specific;
    \item available zone-local controls such as $Q_{AC,i}$ remain zone-specific;
    \item selected non-HVAC disturbances may come from an upstream aggregate
          source.
\end{itemize}

Therefore DEP2 is not a different RC equation.  It is:
\[
\boxed{
\text{DEP1-type physical graph}
+
\text{aggregate-to-zone disturbance allocation operator}.
}
\]

\section*{15. Generic DEP2 aggregate-to-zone allocation}

For disturbance family $g$, let the upstream Phase-B aggregation semantics be
\[
\boxed{
\bar q_g
=
\sum_{i=1}^{N_z} w_{g,i} q_{g,i},
}
\tag{E0.3-13}
\]
where the coefficients $w_{g,i}$ must come from authoritative aggregation
metadata.  They are not inferred from the number of zones.

Define the effective reconstructed zone disturbance:
\[
\boxed{
q_{g,i}^{DEP2}
=
\lambda_{g,i}\bar q_g.
}
\tag{E0.3-14}
\]

To preserve the upstream aggregate exactly,
\[
\boxed{
\sum_{i=1}^{N_z}
w_{g,i}\lambda_{g,i}
=
1.
}
\tag{E0.3-15}
\]

For positive aggregation coefficients, a convenient unconstrained
parameterization is
\[
\boxed{
\lambda_{g,i}
=
\frac{\exp(\alpha_{g,i})}
{\sum_{j=1}^{N_z} w_{g,j}\exp(\alpha_{g,j})}.
}
\tag{E0.3-16}
\]

This guarantees
\[
\lambda_{g,i}>0
\]
and (E0.3-15).

If
\[
\sum_i w_{g,i}=1
\]
and every $\alpha_{g,i}=0$, then
\[
\lambda_{g,i}=1.
\]

For the paper's equal two-zone mean,
\[
w_D=w_K=\frac12,
\]
and (E0.3-15) becomes
\[
\lambda_D+\lambda_K=2,
\]
recovering the controlled PINODE/EPSR case as a special instance.

If the Phase-B aggregation semantics for a source are not expressible by a
known linear coefficient vector, a generic DEP2 allocator must not be silently
constructed.  The allocation contract must instead be supplied explicitly.

\section*{16. Zone total-energy balance}

For zone $i$, define the stored thermal energy relative to an arbitrary fixed
reference temperature $T_{\mathrm{ref}}$:
\[
\boxed{
E_i
=
\sum_{v\in\Vc_i}
C_v(T_v-T_{\mathrm{ref}}).
}
\tag{E0.3-17}
\]

Then
\[
\dot E_i
=
\sum_{v\in\Vc_i}C_v\dot T_v.
\]

All resistance edges whose two endpoints lie inside the same zone cancel
exactly.

Therefore
\[
\boxed{
\begin{aligned}
\dot E_i
={}&
\sum_{\substack{e=\{v,b\}\\
v\in\Vc_i,\ b\in\Vb}}
\frac{T_b-T_v}{R_e}
\\
&+
\sum_{\substack{e=\{v,w\}\\
v\in\Vc_i,\ \zeta(w)\ne i}}
\frac{T_w-T_v}{R_e}
\\
&+
\sum_{s\in\Qset}
\left(
\sum_{v\in\Vc_i}\Gamma_{vs}
\right)q_s.
\end{aligned}
}
\tag{E0.3-18}
\]

This is the generic zone energy-balance identity later consumed by
physics-residual methods and EBP-PINODE.

\section*{17. Building/network total-energy balance}

Define
\[
E_{\mathrm{tot}}
=
\sum_{i\in\Zset}E_i.
\]

Summing (E0.3-18) over all modeled zones cancels every internal and inter-zone
resistance edge.

Hence
\[
\boxed{
\begin{aligned}
\dot E_{\mathrm{tot}}
={}&
\sum_{\substack{e=\{v,b\}\\
v\in\Vc,\ b\in\Vb}}
\frac{T_b-T_v}{R_e}
\\
&+
\sum_{s\in\Qset}
\left(
\sum_{v\in\Vc}\Gamma_{vs}
\right)q_s.
\end{aligned}
}
\tag{E0.3-19}
\]

For conservative routing of source $s$,
\[
\sum_v\Gamma_{vs}=1.
\]

Thus each physical heat source appears exactly once in the global
energy balance.

\section*{18. Parameter contract}

The generic physical parameter set is
\[
\boxed{
\theta_{RC}
=
\left\{
C_v:v\in\Vc
\right\}
\cup
\left\{
R_e:e\in\Er
\right\}.
}
\]

Optional physical routing/allocation parameters are kept separately:
\[
\theta_{\Gamma},
\qquad
\theta_{\Lambda}.
\]

For differentiable estimation:
\[
\boxed{
R_e
=
R_{e,\min}
+
\operatorname{softplus}(\rho_{R_e}),
}
\tag{E0.3-20}
\]
\[
\boxed{
C_v
=
C_{v,\min}
+
\operatorname{softplus}(\rho_{C_v}).
}
\tag{E0.3-21}
\]

When a bounded scalar fraction is used:
\[
\boxed{
\eta=\sigma(\rho_\eta).
}
\tag{E0.3-22}
\]

$R$, $C$, routing fractions, and DEP2 allocation factors are
\emph{model/physical parameters}.  They are not generic hyperparameters.

Hyperparameters such as network width, optimizer settings, numerical rollout
length, or method-specific loss weights belong to later family/method
mathematical contracts.

\section*{19. ScaleBridge canonical graph templates}

The names 1R1C, 2R2C, 3R2C, and 4R3C identify graph signatures.
Heat routing remains an explicit, separate contract.

\subsection*{19.1 1R1C}

State nodes:
\[
\Vc=\{a\}.
\]

Boundary nodes:
\[
\Vb=\{o\}.
\]

Resistance edges:
\[
\Er=\{\{a,o\}\}.
\]

Parameters:
\[
C_a,\qquad R_{ao}.
\]

With the conservative one-node heat routing,
\[
\boxed{
C_a\dot T_a
=
\frac{T_o-T_a}{R_{ao}}
+
Q_c+Q_r.
}
\tag{E0.3-23}
\]

The paper-compatible optional reduced-order participation form is
\[
C_a\dot T_a
=
\frac{T_o-T_a}{R_{ao}}
+
Q_c+\eta_r^{(1C)}Q_r.
\]

\subsection*{19.2 2R2C}

State nodes:
\[
\Vc=\{a,m\}.
\]

Boundary:
\[
\Vb=\{o\}.
\]

Resistance edges:
\[
\Er=
\left\{
\{a,o\},
\{a,m\}
\right\}.
\]

Parameters:
\[
C_a,\ C_m,\ R_{ao},\ R_{am}.
\]

Under the paper-compatible air/mass radiative split:
\[
\boxed{
C_a\dot T_a
=
\frac{T_o-T_a}{R_{ao}}
+
\frac{T_m-T_a}{R_{am}}
+
Q_c
+
(1-\eta_r)Q_r,
}
\tag{E0.3-24}
\]
\[
\boxed{
C_m\dot T_m
=
\frac{T_a-T_m}{R_{am}}
+
\eta_rQ_r.
}
\tag{E0.3-25}
\]

Adding:
\[
\boxed{
C_a\dot T_a+C_m\dot T_m
=
\frac{T_o-T_a}{R_{ao}}+Q_c+Q_r.
}
\tag{E0.3-26}
\]

\subsection*{19.3 3R2C}

ScaleBridge fixes the 3R2C graph signature as two capacitances with three
resistive paths.

State nodes:
\[
\Vc=\{a,m\}.
\]

Boundary:
\[
\Vb=\{o\}.
\]

Resistance edges:
\[
\boxed{
\Er=
\left\{
\{a,o\},
\{a,m\},
\{m,o\}
\right\}.
}
\]

Parameters:
\[
C_a,\ C_m,\ R_{ao},\ R_{am},\ R_{mo}.
\]

For arbitrary routed heat powers $Q_a,Q_m$ satisfying the selected routing
contract:
\[
\boxed{
C_a\dot T_a
=
\frac{T_o-T_a}{R_{ao}}
+
\frac{T_m-T_a}{R_{am}}
+
Q_a,
}
\tag{E0.3-27}
\]
\[
\boxed{
C_m\dot T_m
=
\frac{T_o-T_m}{R_{mo}}
+
\frac{T_a-T_m}{R_{am}}
+
Q_m.
}
\tag{E0.3-28}
\]

The zone total is
\[
\boxed{
C_a\dot T_a+C_m\dot T_m
=
\frac{T_o-T_a}{R_{ao}}
+
\frac{T_o-T_m}{R_{mo}}
+
Q_a+Q_m.
}
\tag{E0.3-29}
\]

No special 3R2C heat split is implied by the topology name; it must be
specified by $\Gamma$.

\subsection*{19.4 4R3C}

ScaleBridge fixes the 4R3C graph signature as zone air, envelope, and internal
mass.

State nodes:
\[
\Vc=\{a,e,m\}.
\]

Boundary:
\[
\Vb=\{o\}.
\]

Resistance edges:
\[
\boxed{
\Er=
\left\{
\{a,o\},
\{a,e\},
\{e,o\},
\{a,m\}
\right\}.
}
\]

Parameters:
\[
C_a,\ C_e,\ C_m,\ R_{ao},\ R_{ae},\ R_{eo},\ R_{am}.
\]

For routed node heat powers $Q_a,Q_e,Q_m$:
\[
\boxed{
C_a\dot T_a
=
\frac{T_o-T_a}{R_{ao}}
+
\frac{T_e-T_a}{R_{ae}}
+
\frac{T_m-T_a}{R_{am}}
+
Q_a,
}
\tag{E0.3-30}
\]
\[
\boxed{
C_e\dot T_e
=
\frac{T_o-T_e}{R_{eo}}
+
\frac{T_a-T_e}{R_{ae}}
+
Q_e,
}
\tag{E0.3-31}
\]
\[
\boxed{
C_m\dot T_m
=
\frac{T_a-T_m}{R_{am}}
+
Q_m.
}
\tag{E0.3-32}
\]

The total balance is
\[
\boxed{
\begin{aligned}
C_a\dot T_a+C_e\dot T_e+C_m\dot T_m
={}&
\frac{T_o-T_a}{R_{ao}}
+
\frac{T_o-T_e}{R_{eo}}\\
&+
Q_a+Q_e+Q_m.
\end{aligned}
}
\tag{E0.3-33}
\]

Again, the graph signature does not silently prescribe heat routing.

\section*{20. Multi-zone construction from elemental zone graphs}

For each modeled zone $i$, instantiate a chosen elemental RC graph
\[
\mathcal{G}_{RC,i}.
\]

The uncoupled network is the disjoint union
\[
\mathcal{G}_{0}
=
\bigsqcup_{i\in\Zset}
\mathcal{G}_{RC,i}.
\]

DEP1/DEP2 physical coupling is produced by adding an explicit set of
inter-zone resistance edges:
\[
\mathcal{E}_{Z}.
\]

Therefore
\[
\boxed{
\mathcal{G}_{multi}
=
\left(
\bigsqcup_{i\in\Zset}
\mathcal{G}_{RC,i}
\right)
\cup
\mathcal{E}_{Z}.
}
\tag{E0.3-34}
\]

Every edge in $\mathcal{E}_{Z}$ is undirected and has one shared positive
resistance.

The construction supports arbitrary $N_z$ and arbitrary allowed thermal
adjacency.  The paper's one Dining--Kitchen air-to-air edge is only the
$N_z=2$ special case.

Detailed adjacency ownership and latent-state initialization are finalized in
E0-4.

\section*{21. Phase B / Phase C / Phase D provider boundary}

The mathematical RC model is written in terms of named ports and does not
depend on filesystem paths or a particular training pipeline.

\subsection*{Training / offline estimation}

The authoritative provider is Phase D through E0-2:
\begin{itemize}
    \item observed zone temperature,
    \item prescribed outdoor temperature,
    \item QAC,
    \item grouped or component heat disturbances,
    \item IND/DEP1/DEP2 information architecture,
    \item aggregation/Phase-C lineage,
    \item temporal partitions.
\end{itemize}

\subsection*{Upstream origin}

Phase B is the authoritative source of measured/aggregated temperature
semantics and aggregation definitions/weights.

Phase C is the authoritative source of learned heat-input quantities such as
QAC, QSol1/QSol2, and internal-gain components that Phase D subsequently
materializes/groups.

\subsection*{Runtime}

Later runtime composition may supply the same thermal ports from live/replayed
Phase-B measurements and Phase-C model runtimes.

The RC equations remain unchanged.

PHVAC is computed downstream from the corresponding QAC trajectory and does
not enter the RC balance.

\section*{22. Generic physical invariants required of every backend}

Any NumPy, PyTorch, Neuromancer, CasADi, Pyomo, or Bayesian representation
must satisfy the same contract.

Required invariants include:

\begin{enumerate}
    \item $R_e>0$ for every resistance.
    \item $C_v>0$ for every capacitance.
    \item Internal edge flow is antisymmetric.
    \item A shared inter-zone edge has one parameter, not two directional
          parameters.
    \item Internal edges cancel from the appropriate zone/building energy sum.
    \item Conservative heat routing satisfies (E0.3-6).
    \item DEP2 allocation satisfies (E0.3-15).
    \item PHVAC never appears in the thermal RHS.
    \item E0-2 exact signal bindings are respected; no substring inference.
    \item A missing applicable signal is not silently zero-filled.
    \item The observed-state operator selects the exact Phase-D zone-air states.
    \item Latent states remain physical topology states rather than arbitrary
          learned coordinates.
    \item Named RC templates reproduce the generic graph equation exactly.
\end{enumerate}

\section*{23. Separation from later Phase E steps}

This contract intentionally stops at the continuous-time physical model.

E0-4 owns:
\begin{itemize}
    \item arbitrary-zone adjacency/coupling realization,
    \item detailed shared inter-zone resistance registry,
    \item latent-state initialization,
    \item aggregation-aware DEP2 allocator binding to actual Phase-B metadata.
\end{itemize}

E0-5 owns:
\begin{itemize}
    \item Euler,
    \item Heun/RK2,
    \item RK4,
    \item the precisely identified sixth-order RK scheme,
    \item physical/normalized time,
    \item substep forcing interpolation/ZOH.
\end{itemize}

E0-6 owns backend parity.

E0-7 owns runtime composition and PHVAC runtime.

E0-8 owns Optuna/MLflow lifecycle.

Specific ML, SciML, Optimization, and Bayesian methods must still receive
their own detailed versioned mathematical \texttt{.tex} contracts before code
is written.

\section*{24. Lock summary}

The proposed E0-3 mathematical lock is:

\begin{enumerate}
    \item One generic graph equation governs every RC topology.
    \item 1R1C/2R2C/3R2C/4R3C are graph templates, not duplicated solvers.
    \item Heat routing is explicit and separate from graph topology.
    \item 1C/nC is a state-structure distinction.
    \item All-to-One/IND/DEP1/DEP2 is an information-architecture distinction.
    \item DEP1 and DEP2 may share the same physical graph.
    \item DEP2 differs by an explicit aggregate-to-zone information allocator.
    \item The generic DEP2 conservation rule is
          $\sum_i w_{g,i}\lambda_{g,i}=1$.
    \item Phase D/E0-2 supplies authoritative training-time signal bindings.
    \item Phase B/C semantics remain upstream; the RC core consumes named ports.
    \item QAC is thermal control; PHVAC is downstream electrical output.
    \item State temperatures are physically named; latent states are topology
          states.
    \item All $R/C$ parameters remain positive.
    \item Total-energy identities are generated from the graph and must hold
          independently of method/backend.
\end{enumerate}

\end{document}
